\documentclass[11pt,a4paper]{article}
\usepackage[utf8]{inputenc}
\usepackage[T1]{fontenc}
\usepackage{lmodern}
\usepackage{amsmath,amssymb,amsthm}
\usepackage[margin=2.6cm]{geometry}
\usepackage{booktabs}
\usepackage{listings}
\usepackage[hidelinks]{hyperref}
\hypersetup{pdftitle={Erdos Problem 700(a) for two prime factors and small-order checks of open Graffiti.pc conjectures},pdfauthor={Verifiable Mathematical Reasoning Corpus contribution}}

\newtheorem{theorem}{Theorem}[section]
\newtheorem{lemma}[theorem]{Lemma}
\newtheorem{proposition}[theorem]{Proposition}
\newtheorem{corollary}[theorem]{Corollary}
\newtheorem{conjecture}[theorem]{Conjecture}
\theoremstyle{definition}
\newtheorem{definition}[theorem]{Definition}
\newtheorem{remark}[theorem]{Remark}

\newcommand{\res}{\operatorname{res}}
\newcommand{\im}{\operatorname{im}}
\newcommand{\rad}{\operatorname{rad}}
\newcommand{\diam}{\operatorname{diam}}
\newcommand{\tr}{\operatorname{tree}}

\lstset{basicstyle=\ttfamily\scriptsize,breaklines=true,columns=fullflexible,
  frame=single,framerule=0.3pt,keepspaces=true,showstringspaces=false}

\title{Erd\H{o}s Problem \#700(a) for integers with two prime factors,\\
and exhaustive small-order checks of eight open Graffiti.pc conjectures}
\author{Contribution to the \emph{Verifiable Mathematical Reasoning Corpus}\\
\small (prepared with AI assistance; all computations reproducible from the appendix)}
\date{November 2026}

\begin{document}
\maketitle

\begin{abstract}
For composite $n$ let $f(n)=\min_{1<k\le n/2}\gcd\bigl(n,\binom nk\bigr)$ and let $P(n)$ be
the largest prime factor of $n$. Erd\H{o}s and Szekeres asked to characterize the composite
$n$ with $f(n)=n/P(n)$ (Erd\H{o}s Problem \#700(a)). This question is still open; on the
problem's forum it has been described as probably admitting no satisfactory answer in
general. We prove three things.
\begin{itemize}
\item A general necessary condition: if $n$ is not a prime power, then $P(n)$ divides $n$
  exactly, and every other prime-power component $p^{v_p(n)}$ is smaller than $P(n)$.
\item A complete characterization for every composite $n$ with at most two distinct prime
  factors (Theorem~\ref{thm:A}). For $n=p^aq^b$ with $p<q$, equality holds if and only if
  $b=1$, $p^a<q$, and an explicit base-$p$ carry condition on small multiples of $q$ is
  satisfied.
\item An infinitude statement: for every prime power $p^a$, all primes $q>p^a$ with
  $q\equiv-1\pmod{p^{2a}}$ satisfy $f(p^aq)=n/P(n)$.
\end{itemize}
The characterization agrees with brute force for all $11\,237$ such $n\le30\,000$.

In a second, purely computational part we evaluate eight Graffiti.pc conjectures from
\emph{Written on the Wall~II} that are open in the Formal Conjectures repository: 19, 40, 61,
72, 96, 133, 189 and 199. We use the repository's exact definitions and check every
connected graph with at most ten vertices; no counterexample exists. We also verify, on the
same graphs, the inequality $f(G)\ge\res(G)+\im(G)$, which was left open in corpus item
[066] and implies WOWII~61. We record three short arguments: a proof of WOWII~96 (a proof
with the same idea was posted earlier by another author), the reduction of WOWII~133 to
$C_4$-free graphs, and WOWII~189 for graphs of diameter at most two.

None of these is a complete resolution of a named problem by the present work; the precise
status of each item is stated where it appears and summarized in
Section~\ref{sec:assessment}.
\end{abstract}

\tableofcontents

%=====================================================================
\section{Erd\H{o}s Problem \#700(a)}
%=====================================================================

\subsection{Background}
Erd\H{o}s and Szekeres~\cite{ErSz78} introduced
$f(n)=\min_{1<k\le n/2}\gcd(n,\binom nk)$ for composite $n$ and asked three questions,
recorded as Erd\H{o}s Problem \#700~\cite{EP700}:
\begin{itemize}
\item[(a)] characterize the composite $n$ with $f(n)=n/P(n)$;
\item[(b)] are there infinitely many composite $n$ with $f(n)>n^{1/2}$?
\item[(c)] is $f(n)\ll_A n/(\log n)^A$ for every $A>0$?
\end{itemize}
Question~(b) has been answered affirmatively in 2026, by products of three primes; see the
problem page and corpus item [086]. Questions (a) and (c) are listed as open (page last
edited 28 August 2026). It is classical that $f(pq)=n/P(n)$ for $n=pq$, and that $f(30)=6$.
The forum discussion records further squarefree examples, such as $2310$ and
$71677177=11\cdot13\cdot19\cdot23\cdot31\cdot37$, and suggests that a general
characterization is hopeless~\cite{EP700}. Our results give the complete answer for $n$
with at most two prime factors.

\subsection{Tools}
For a prime $\ell$ and nonnegative integers $x,y$, let $\kappa_\ell(x,y)$ be the number of
carries when $x$ and $y$ are added in base $\ell$, and let $s_\ell(x)$ be the base-$\ell$
digit sum.

\begin{lemma}[Kummer]\label{lem:kummer}
$v_\ell\binom{x+y}{x}=\kappa_\ell(x,y)=\bigl(s_\ell(x)+s_\ell(y)-s_\ell(x+y)\bigr)/(\ell-1)$.
Moreover $\kappa_\ell(x,y)=0$ if and only if every base-$\ell$ digit of $x$ is at most the
corresponding digit of $x+y$.
\end{lemma}

Consequently
\begin{equation}\label{eq:gcd}
\gcd\Bigl(n,\binom nk\Bigr)=\prod_{\ell\mid n}\ell^{\min\{v_\ell(n),\,\kappa_\ell(k,n-k)\}}.
\end{equation}

\begin{lemma}\label{lem:lower}
For $1\le k\le n$ and every prime $\ell$,
$v_\ell\binom nk\ge v_\ell(n)-v_\ell(k)$.
\end{lemma}
\begin{proof}
$k\binom nk=n\binom{n-1}{k-1}$.
\end{proof}

\begin{lemma}\label{lem:component}
Let $\ell^e\,\|\,n$ with $e\ge1$, and suppose $\ell^e\le n/2$. Then
$\gcd(n,\binom{n}{\ell^e})=n/\ell^e$.
\end{lemma}
\begin{proof}
In base $\ell$ the lowest nonzero digit of $n$ sits at position $e$. Adding $\ell^e$ to
$n-\ell^e$ restores that digit without a carry, so $v_\ell\binom n{\ell^e}=0$. For every
other prime $r\mid n$, Lemma~\ref{lem:lower} gives $v_r\binom n{\ell^e}\ge v_r(n)$.
\end{proof}

\subsection{A necessary condition for general $n$}
\begin{proposition}\label{prop:necessary}
Let $n$ be composite and not a prime power, with $P=P(n)$. If $f(n)=n/P$, then
$P\,\|\,n$, and $p^{v_p(n)}<P$ for every prime $p\mid n$ with $p\ne P$.
\end{proposition}
\begin{proof}
Every prime-power component $\ell^e\,\|\,n$ satisfies $\ell^e\le n/2$, so
Lemma~\ref{lem:component} gives $f(n)\le n/\ell^e$. With $\ell=P$ this forces $e=1$. With
$\ell=p\ne P$ it gives $n/P\le n/p^{v_p(n)}$, i.e.\ $p^{v_p(n)}\le P$, and equality is
impossible for distinct primes.
\end{proof}

For prime powers, $f(p^a)=p$ for all $a\ge2$ (Theorem~\ref{thm:A}(i)). So among prime powers,
$f(n)=n/P(n)$ holds exactly for the squares $p^2$.

\subsection{The case of at most two prime factors}
\begin{theorem}\label{thm:A}
\begin{enumerate}
\item[(i)] If $n=p^a$ with $a\ge2$, then $f(n)=p$. Hence $f(n)=n/P(n)$ if and only if
  $a=2$.
\item[(ii)] Let $n=p^aq^b$ with primes $p<q$ and $a,b\ge1$, so that $P(n)=q$. Then
  $f(n)=n/q$ if and only if all of the following hold:
  \begin{enumerate}
  \item[(1)] $b=1$;
  \item[(2)] $p^a<q$;
  \item[(3)] for every $t\in\{1,\dots,a-1\}$ and every integer $j'\ge1$ with $p\nmid j'$
    and $p^tj'\le p^a/2$,
    \[
    \kappa_p\bigl(j'q,\ (p^{a-t}-j')q\bigr)\ \ge\ a .
    \]
  \end{enumerate}
\end{enumerate}
\end{theorem}

\begin{proof}
(i) For $2\le k\le p^{a}/2$ we have $v_p(k)\le a-1$, so Lemma~\ref{lem:lower} gives
$p\mid\gcd(p^a,\binom{p^a}k)$. At $k=p^{a-1}$, adding $p^{a-1}$ and $(p-1)p^{a-1}$ in base
$p$ produces exactly one carry, so the gcd equals $p$.

(ii) The values $k=q^b$ and $k=p^a$ are at most $n/2$, so Lemma~\ref{lem:component} gives
$f(n)\le n/q^b$ and $f(n)\le n/p^a$. If $b\ge2$, then $f(n)<n/q$. If $p^a>q$, then
$f(n)\le n/p^a<n/q$. Hence (1) and (2) are necessary.

Assume (1) and (2). Then $f(n)\le\gcd(n,\binom nq)=p^a$, so $f(n)=p^a$ if and only if
$\gcd(n,\binom nk)\ge p^a$ for every $2\le k\le n/2$. By \eqref{eq:gcd} this gcd equals
$p^{\min\{a,\kappa_p\}}q^{\min\{1,\kappa_q\}}$. If $\kappa_q(k,n-k)\ge1$, the gcd is at
least $q>p^a$.

Suppose $\kappa_q(k,n-k)=0$. Since $p^a<q$, the base-$q$ digits of $n=p^aq$ are $(p^a,0)$.
By Lemma~\ref{lem:kummer}, $k=jq$ with $0\le j\le p^a$, and $2\le k\le n/2$ means
$1\le j\le p^a/2$. For such $k$ the gcd is $p^{\min\{a,\kappa_p(jq,(p^a-j)q)\}}$.

Write $j=p^tj'$ with $p\nmid j'$; then $t\le a-1$. Multiplying both summands by $p^t$ only
shifts digits, so $\kappa_p(jq,(p^a-j)q)=\kappa_p(j'q,(p^{a-t}-j')q)$. If $t=0$,
Lemma~\ref{lem:lower} applied to $\binom{p^aq}{jq}$ gives at least $a$ carries
automatically. Hence the condition reduces exactly to (3).
\end{proof}

\begin{corollary}\label{cor:small}
Let $q$ be prime.
\begin{enumerate}
\item $f(pq)=n/P(n)$ for $p<q$ (classical: condition (3) is vacuous).
\item $f(4q)=4=n/P(n)$ for every prime $q\ge5$.
\item $f(9q)=9=n/P(n)$ if and only if $q\ge11$ and $s_3(2q)\ge4$.
\item $f(8q)=8=n/P(n)$ if and only if $q\ge11$, $s_2(q)\ge3$ and $s_2(3q)\ge3$.
\item For $n=p^2q$ with $p^2<q$: $f(n)=n/P(n)$ if and only if
  $\kappa_p(j'q,(p-j')q)\ge2$ for all $1\le j'\le p/2$.
\end{enumerate}
\end{corollary}
\begin{proof}
Specialize (3). For $p=2,a=2$ the only case is $t=1$, $j'=1$, with
$\kappa_2(q,q)=s_2(q)\ge2$. For $p=3,a=2$: $\kappa_3(q,2q)=s_3(2q)/2$. For $p=2,a=3$ the
cases are $(t,j')=(1,1)$, giving $\kappa_2(q,3q)=s_2(3q)$, and $(2,1)$, giving $s_2(q)$.
\end{proof}

For example $f(104)=8$ ($q=13$), while $f(88)<8$ ($s_2(33)=2$) and $f(136)<8$
($s_2(17)=2$).

\begin{corollary}[Infinitude]\label{cor:infinite}
Let $p$ be prime and $a\ge1$. Every prime $q>p^a$ with $q\equiv-1\pmod{p^{2a}}$ satisfies
$f(p^aq)=p^a=n/P(n)$. In particular, for every prime power $p^a$ there are infinitely many
primes $q$ with $f(p^aq)=n/P(n)$; this includes integers with arbitrarily large non-squarefree
part.
\end{corollary}
\begin{proof}
Fix $t\in[1,a-1]$ and an admissible $j'$, and put $x=j'q$, $y=(p^{a-t}-j')q$. A carry leaves
position $i-1$ exactly when $(x\bmod p^i)+(y\bmod p^i)\ge p^i$. For $a-t\le i\le 2a$ we
have $j'<p^{a-t}\le p^i$ and $0<p^{a-t}-j'<p^i$. Since $q\equiv-1\pmod{p^i}$, this gives
$x\bmod p^i=p^i-j'$ and $y\bmod p^i=p^i-p^{a-t}+j'$, with sum $2p^i-p^{a-t}\ge p^i$. That is
$2a-(a-t)+1\ge a$ carries, so condition (3) holds. Conditions (1) and (2) hold by
assumption, and Dirichlet's theorem supplies infinitely many such primes.
\end{proof}

\subsection{Verification and data}
A direct C program computes $f(n)$ from \eqref{eq:gcd}. Together with a Python
implementation of Theorem~\ref{thm:A} (Appendix~\ref{app:code}), it confirms the theorem for
all $11\,237$ composite $n\le30\,000$ with at most two prime factors, with no mismatch.
Outside the semiprimes $pq$, there are $7\,722$ composite $n\le30\,000$ with
$f(n)=n/P(n)$.

For question~(c), primorials are a natural test family. Here $f(n)$ falls rapidly below
$n/P(n)$, consistent with a positive answer to (c):

\begin{center}
\begin{tabular}{rrrr}
\toprule
$n$ & $f(n)$ & $n/P(n)$ & $f(n)/(n/P(n))$\\
\midrule
$30$ & $6$ & $6$ & $1$\\
$210$ & $14$ & $30$ & $0.467$\\
$2310$ & $154$ & $210$ & $0.733$\\
$30030$ & $165$ & $2310$ & $0.0714$\\
$510510$ & $510$ & $30030$ & $0.0170$\\
$9699690$ & $2730$ & $510510$ & $0.0053$\\
\bottomrule
\end{tabular}
\end{center}

\paragraph{Status.} Theorem~\ref{thm:A} settles question (a) completely in the class
$\omega(n)\le2$. Corpus item [086] settles certain products of three close primes. The
general question (a) remains open.

%=====================================================================
\section{Eight open Graffiti.pc conjectures}
%=====================================================================

\subsection{Definitions}
We use the formalizations in the Formal Conjectures repository~\cite{FC}, where these eight
statements are tagged ``research open'' as of 6 November 2026:
\begin{itemize}
\item $l(v)$: the independence number of $G[N(v)]$, with $\Lambda=\max_vl(v)$ and
  $\bar l=\frac1n\sum_vl(v)$;
\item $b$: the largest order of an induced bipartite subgraph;
\item $f$: the largest order of an induced forest;
\item $\tr$: the largest order of an induced tree;
\item $\mathrm{path}$: the largest order of an induced path;
\item $p$: the path-cover number;
\item $\res$: the Havel--Hakimi residue;
\item $\sigma$: the second-smallest degree;
\item $\kappa$: the vertex connectivity;
\item $\mathrm{distEven}(v)$: the number of vertices at even distance from $v$, counting
  $v$ itself.
\end{itemize}

The conjectures checked are:
\begin{itemize}
\item[19:] $b\ge\lfloor\bar e+\Lambda\rfloor$, where $\bar e$ is the average eccentricity;
\item[40:] $f\ge\lceil(p+b+1)/2\rceil$;
\item[61:] $f\ge\res+\lceil\diam/3\rceil$;
\item[72:] $\tr\ge\lceil\bar e+\Lambda/3\rceil$;
\item[96:] $\alpha\le1+\min_v\mathrm{distEven}(v)\,(\Lambda-1)$;
\item[133:] $\mathrm{path}\ge\rad+\bar l^{\,cC_4}$, where $cC_4=0$ if $G$ contains a
  (not necessarily induced) $4$-cycle and $cC_4=1$ otherwise;
\item[189:] if $\max_v\mathrm{distEven}(v)\le1+\sigma$, then $G$ has a Hamiltonian path;
\item[199:] if $\tr-2\le\kappa$, then $G$ has a Hamiltonian path.
\end{itemize}

\subsection{Exhaustive verification up to ten vertices}
\begin{theorem}[computational]\label{thm:wow}
None of the eight conjectures above has a counterexample among the connected graphs with at
most $10$ vertices.
\end{theorem}
The graphs were generated by \texttt{nauty-geng -c}; the counts of connected graphs for
$n=2,\dots,10$ are $1,2,6,21,112,853,11117,261080,11716571$. Every invariant is computed
exactly by subset enumeration or dynamic programming (program \texttt{wow.c},
Appendix~\ref{app:code}). The program was cross-checked against independent
\texttt{networkx} brute force on test graphs, including the Petersen graph.

\begin{remark}
In October 2026 several of these statements received public claims on the Formal
Conjectures issue tracker:
\begin{itemize}
\item a counterexample to 72, with $31$ vertices (issue~\#7662; an earlier record is
  reported there);
\item proofs of 19 (pull requests \#4559, \#7571) and of 96 (issue~\#7664).
\end{itemize}
By Theorem~\ref{thm:wow}, every counterexample to 72 has at least $11$ vertices. We found no
public claims for 61, 133, 189 or 199.
\end{remark}

\subsection{The residue--induced-matching inequality}
Corpus item [066] reduced WOWII~61 to the following inequality and left it open:
\begin{equation}\label{eq:cand}
f(G)\ \ge\ \res(G)+\im(G),
\end{equation}
where $\im$ is the induced matching number. The reduction runs as follows. The edges with
indices $0,3,6,\dots$ of a diametral geodesic form an induced matching of size
$\lceil\diam/3\rceil$.

\begin{proposition}[computational]
Inequality \eqref{eq:cand} holds for every graph with at most $10$ vertices.
\end{proposition}
\begin{proof}
Connected graphs were checked exhaustively (program \texttt{l2.c}). For disconnected graphs,
$f$ and $\im$ are additive over components and $\res$ is subadditive (corpus [066],
Lemma ``disjoint-union subadditivity''), so the connected case suffices.
\end{proof}

The variant with $f$ replaced by $\alpha_2$, the largest order of an induced subgraph of
maximum degree at most one, already fails at $5$ vertices.

\subsection{Three short arguments}

\begin{theorem}[WOWII 96]\label{thm:96}
For every connected graph $G$ with at least two vertices,
$\alpha(G)\le1+\min_v\mathrm{distEven}(v)\,(\Lambda-1)$.
\end{theorem}
\begin{proof}
If $\Lambda=1$, every neighbourhood is a clique, so the connected graph $G$ is complete and
$\alpha=1$. Assume $\Lambda\ge2$.

Fix $v$, put $D=\mathrm{distEven}(v)$, and let $L_i$ be the set of vertices at distance $i$
from $v$. Assign to each $x\in L_{2j+1}$ a neighbour $\pi(x)\in L_{2j}$. For $y$ at even
distance let $C(y)=\pi^{-1}(y)\subseteq N(y)\cap L_{2j+1}$. Then $V$ is the disjoint union
of the sets $\{y\}\cup C(y)$.

Let $I$ be independent. If $y\in I$, then $I\cap C(y)=\emptyset$. If $y=v\notin I$, then
$|I\cap C(v)|\le l(v)\le\Lambda$. If $y\in L_{2j}$ with $j\ge1$ and $y\notin I$, choose a
neighbour $z\in L_{2j-1}$ of $y$. This $z$ has no neighbour in $L_{2j+1}$, so
$(I\cap C(y))\cup\{z\}$ is independent in $G[N(y)]$, and $|I\cap C(y)|\le l(y)-1\le\Lambda-1$.

Summing over the $D$ even-distance vertices gives
$|I|\le\Lambda+(D-1)(\Lambda-1)=1+D(\Lambda-1)$.
\end{proof}
\noindent\emph{Priority.} A proof by the same BFS-layer idea was posted for review by
J.~J.~Gebendorfer on 28 October 2026 (Formal Conjectures issue~\#7664). We include our
argument only for completeness and claim no priority.

\begin{proposition}[WOWII 133, $C_4$ branch]
If $G$ is connected and contains a $4$-cycle, then $\mathrm{path}(G)\ge\rad(G)+1$. Hence
WOWII~133 is equivalent to its restriction to $C_4$-free graphs. In a $C_4$-free graph
$l(v)=d(v)-t(v)$, where $t(v)$ is the number of edges inside $N(v)$.
\end{proposition}
\begin{proof}
A diametral geodesic is an induced path on $\diam+1\ge\rad+1$ vertices. In a $C_4$-free
graph no vertex of $N(v)$ has two neighbours in $N(v)$, so $G[N(v)]$ is a matching plus
isolated vertices.
\end{proof}

\begin{proposition}[WOWII 189 for diameter $\le2$]
If $\diam(G)\le2$ and $\max_v\mathrm{distEven}(v)\le1+\sigma(G)$, then $G$ has a Hamiltonian
path.
\end{proposition}
\begin{proof}
For diameter at most two, $\mathrm{distEven}(v)=n-d(v)$, so the hypothesis reads
$\delta+\sigma\ge n-1$. Any two distinct vertices therefore have degree sum at least
$n-1$, and Ore's theorem for Hamiltonian paths applies.
\end{proof}

%=====================================================================
\section{A status note on order-five Austin laws}
%=====================================================================
The Equational Theories Project lists $96$ order-five laws whose finite models are trivial
and whose infinite models were unknown~\cite{ETPblueprint}. Since then, classifications have
come from three sources:
\begin{itemize}
\item Zhao--Wu--Yang--Miao~\cite{ZWYM} certified $32$ of them as Austin laws, $4$ of which
  overlap the ALPS benchmark positives;
\item ALPS~\cite{ALPS} has two further positives, $12857/33436$;
\item corpus item [085] settled $5093/41179$ and $5107/40951$.
\end{itemize}
In the same-resource baseline of~\cite{ZWYM}, Twee proved $x\approx y$ from $5834$ and
from $40037$, so this dual pair is trivial rather than Austin. Subtracting all of these
leaves $28$ dual pairs, i.e.\ $56$ laws, that are unclassified as far as we can determine:
\begin{center}\small
\begin{tabular}{llll}
5833/40070 & 5837/40221 & 5947/40057 & 5951/40208\\
6820/39485 & 6878/39126 & 6895/39163 & 6912/39214\\
7587/38316 & 7763/38565 & 8485/37519 & 9337/36867\\
9603/36514 & 9663/36487 & 10218/35685 & 10222/35836\\
11082/34889 & 11205/35100 & 11280/34778 & 12073/33998\\
12087/33884 & 12234/33883 & 12883/33020 & 13764/32294\\
13849/32281 & 13992/32280 & 18137/27863 & 18212/27859\\
\end{tabular}
\end{center}
This is bookkeeping, not a theorem. We attempted no classification, because no working
equational prover was available in our environment.

%=====================================================================
\section{Assessment}\label{sec:assessment}
%=====================================================================
\begin{itemize}
\item \textbf{Erd\H{o}s \#700(a).} Partial answer. Proposition~\ref{prop:necessary} and
  Theorem~\ref{thm:A} with Corollaries~\ref{cor:small}--\ref{cor:infinite} are new as far as
  we can determine, and they completely settle the class $\omega(n)\le2$. The general
  question remains open, and the problem's own forum expects no simple answer.
\item \textbf{WOWII.} Only computational evidence, plus three short arguments. Of these,
  Theorem~\ref{thm:96} has an earlier public claim, and the other two are partial. No
  WOWII conjecture is newly resolved here.
\item \textbf{Austin laws.} A status list only.
\end{itemize}
Measured against the corpus criteria (a pre-existing named problem, unchanged and
unconditionally and completely resolved), none of the items is a complete resolution. The
closest is the complete two-prime characterization for Erd\H{o}s \#700(a).

\begin{thebibliography}{9}
\bibitem{ErSz78} P.~Erd\H{o}s and G.~Szekeres, \emph{Some number theoretic problems on
  binomial coefficients}, Austral.\ Math.\ Soc.\ Gaz.\ \textbf{5} (1978), 97--99.
\bibitem{EP700} T.~F.~Bloom, \emph{Erd\H{o}s Problem \#700},
  \texttt{https://www.erdosproblems.com/700}, and its discussion thread (accessed
  8 November 2026).
\bibitem{FC} The Formal Conjectures Authors, \emph{Formal Conjectures},
  \texttt{github.com/google-deepmind/formal-conjectures}, files
  \texttt{WrittenOnTheWallII/GraphConjecture\{19,40,61,72,96,133,189,199\}.lean}.
\bibitem{DeLaVina} E.~DeLaVi\~na, \emph{Written on the Wall II (Conjectures of
  Graffiti.pc)}, University of Houston--Downtown.
\bibitem{ETPblueprint} Equational Theories Project, \emph{Blueprint}, chapter ``Order 5
  Austin laws'', \texttt{github.com/teorth/equational\_theories}.
\bibitem{ZWYM} J.~Zhao, B.~Wu, T.~Yang and X.~Miao, \emph{Trace-tree magmas},
  arXiv:2609.05690 (2026); artifacts at \texttt{github.com/YanbiaoLab/trace-tree-magmas}.
\bibitem{ALPS} E.~Xie, W.~Ye and A.~Zhang, \emph{ALPS}, arXiv:2608.15979 (2026).
\end{thebibliography}

\appendix
\section{Programs}\label{app:code}
\subsection*{A.1 \texttt{f700list.c}: composite $n\le N$ with $f(n)=n/P(n)$}
\begin{lstlisting}[language=C]
#include <stdio.h>
#include <stdlib.h>
typedef unsigned long long u64;
static int carries(u64 a,u64 b,u64 p,int cap){ int c=0,carry=0; while(a||b||carry){ u64 s=a%p+b%p+carry; carry=s>=p; if(carry){c++; if(c>=cap) return cap;} a/=p; b/=p;} return c; }
int main(int argc,char**argv){ u64 N=strtoull(argv[1],0,10);
  for(u64 n=4;n<=N;n++){
    u64 pr[20]; int ex[20],np=0; u64 m=n; for(u64 p=2;p*p<=m;p++) if(m%p==0){pr[np]=p;ex[np]=0;while(m%p==0){m/=p;ex[np]++;}np++;} if(m>1){pr[np]=m;ex[np]=1;np++;}
    if(np==1 && ex[0]==1) continue; /* prime */
    u64 P=pr[np-1], target=n/P, best=n;
    for(u64 k=2;k<=n/2;k++){ u64 g=1; for(int i=0;i<np;i++){ int c=carries(k,n-k,pr[i],ex[i]); for(int t=0;t<c;t++) g*=pr[i]; if(g>=best) break; } if(g<best){best=g; if(best<target) break;} }
    if(best==target && !(np==2 && ex[0]==1 && ex[1]==1)) { printf("%llu =",n); for(int i=0;i<np;i++) printf(" %llu^%d",pr[i],ex[i]); printf("\n"); }
  }
}
\end{lstlisting}
\subsection*{A.2 \texttt{thmA.py}: Theorem~\ref{thm:A} against brute force}
\begin{lstlisting}[language=Python]
# Verify Theorem A (Erdos #700(a), omega(n)<=2) against brute force f(n) for all composite n<=N.
import sys
from sympy import factorint
N=int(sys.argv[1])
def carries(a,b,p):
    c=0; carry=0
    while a or b or carry:
        s=a%p+b%p+carry; carry=1 if s>=p else 0; c+=carry; a//=p; b//=p
    return c
def f_brute(n,fac):
    best=n
    for k in range(2,n//2+1):
        g=1
        for p,e in fac.items():
            g*=p**min(e,carries(k,n-k,p))
            if g>=best: break
        if g<best: best=g
    return best
def predict(n,fac):
    ps=sorted(fac)
    if len(ps)==1:
        p,a=ps[0],fac[ps[0]]
        return a==2
    if len(ps)==2:
        p,q=ps  # p<q, q=P(n)
        a,b=fac[p],fac[q]
        if b!=1 or p**a>q: return False
        for t in range(1,a):
            for jp in range(1, p**(a-t)//2*2+1):
                if jp%p==0: continue
                if p**t*jp > p**a//2: break
                if carries(jp*q,(p**(a-t)-jp)*q,p) < a: return False
        return True
    return None
mism=0; cnt=0; true_count=0
for n in range(4,N+1):
    fac=factorint(n)
    if len(fac)==1 and list(fac.values())[0]==1: continue
    if len(fac)>2: continue
    cnt+=1
    P=max(fac)
    actual = (f_brute(n,fac)==n//P)
    pred=predict(n,fac)
    if actual: true_count+=1
    if pred!=actual:
        mism+=1
        if mism<=10: print("MISMATCH",n,fac,actual,pred)
print(f"checked {cnt} composite n<= {N} with omega<=2; f=n/P holds for {true_count}; mismatches={mism}")
\end{lstlisting}
\subsection*{A.3 \texttt{wow.c}: the eight WOWII conjectures on graph6 input}
\begin{lstlisting}[language=C]
/* Exhaustive checker for open WOWII conjectures 19,40,61,72,96,133,189,199
   (definitions follow the Formal Conjectures Lean formalisation). Input: graph6 lines (connected graphs). */
#include <stdio.h>
#include <stdlib.h>
#include <string.h>
typedef unsigned int U;
static int n; static U adj[16];
static int pc(U x){ return __builtin_popcount(x); }
static int parse(const char*s){
  n=s[0]-63; memset(adj,0,sizeof adj);
  int k=0; const char*p=s+1; int bits=0,val=0;
  for(int j=1;j<n;j++) for(int i=0;i<j;i++){
    if(bits==0){ val=(*p++)-63; bits=6; }
    bits--; if((val>>bits)&1){ adj[i]|=1u<<j; adj[j]|=1u<<i; }
    k++;
  }
  return n;
}
static U reach(U mask,int s){ /* vertices reachable from s inside mask */
  U seen=1u<<s, fr=seen;
  while(fr){ U nf=0; for(U f=fr;f;f&=f-1){ int v=__builtin_ctz(f); nf|=adj[v]&mask; } nf&=~seen; seen|=nf; fr=nf; }
  return seen;
}
static int connected(U mask){ if(!mask) return 0; return reach(mask,__builtin_ctz(mask))==mask; }
static int edges(U mask){ int e=0; for(U m=mask;m;m&=m-1){ int v=__builtin_ctz(m); e+=pc(adj[v]&mask);} return e/2; }
static int ncomp(U mask){ int c=0; while(mask){ U r=reach(mask,__builtin_ctz(mask)); mask&=~r; c++; } return c; }
static int bipartite(U mask){ /* 2-colour each component by BFS */
  U rem=mask;
  while(rem){ int s=__builtin_ctz(rem); U A=1u<<s,B=0, frA=A, frB=0, seen=A;
    while(frA||frB){ U nB=0,nA=0;
      for(U f=frA;f;f&=f-1){ int v=__builtin_ctz(f); nB|=adj[v]&mask; }
      for(U f=frB;f;f&=f-1){ int v=__builtin_ctz(f); nA|=adj[v]&mask; }
      if((nB&A)||(nA&B)) return 0;
      frB=nB&~B; frA=nA&~A; B|=nB; A|=nA; seen|=nA|nB; }
    if(A&B) return 0;
    rem&=~seen; }
  return 1;
}
static int indepnum_in(U mask){ /* max independent set inside mask, simple branching */
  if(!mask) return 0;
  int v=__builtin_ctz(mask);
  /* choose vertex of max degree in mask for branching */
  int best=-1,bd=-1; for(U m=mask;m;m&=m-1){ int u=__builtin_ctz(m); int d=pc(adj[u]&mask); if(d>bd){bd=d;best=u;} }
  if(bd==0) return pc(mask);
  v=best;
  int a=indepnum_in(mask&~(1u<<v));            /* exclude v */
  int b=1+indepnum_in(mask&~(1u<<v)&~adj[v]);   /* include v */
  return a>b?a:b;
}
int main(int argc,char**argv){
  char line[256]; long long cnt=0; long long viol[8]={0}; const int ids[8]={19,40,61,72,96,133,189,199};
  int printed[8]={0};
  static unsigned char hp[1<<12][12]; static int pcov[1<<12];
  while(fgets(line,sizeof line,stdin)){
    line[strcspn(line,"\n")]=0; if(!line[0]) continue;
    parse(line); if(n<2) continue; cnt++;
    U full=(n==32)?~0u:((1u<<n)-1);
    /* distances */
    int dist[16][16];
    for(int s=0;s<n;s++){ for(int t=0;t<n;t++) dist[s][t]=-1; dist[s][s]=0; U seen=1u<<s, fr=seen; int d=0;
      while(fr){ d++; U nf=0; for(U f=fr;f;f&=f-1) nf|=adj[__builtin_ctz(f)]; nf&=~seen; for(U f=nf;f;f&=f-1) dist[s][__builtin_ctz(f)]=d; seen|=nf; fr=nf; } }
    int ecc[16], sumEcc=0, rad=99, diam=0;
    for(int v=0;v<n;v++){ int e=0; for(int t=0;t<n;t++) if(dist[v][t]>e) e=dist[v][t]; ecc[v]=e; sumEcc+=e; if(e<rad) rad=e; if(e>diam) diam=e; }
    int lv[16], maxL=0, sumL=0;
    for(int v=0;v<n;v++){ lv[v]=indepnum_in(adj[v]); sumL+=lv[v]; if(lv[v]>maxL) maxL=lv[v]; }
    int alpha=indepnum_in(full);
    int distEven[16], minDE=99, maxDE=0;
    for(int v=0;v<n;v++){ int c=0; for(int t=0;t<n;t++) if(dist[v][t]%2==0) c++; distEven[v]=c; if(c<minDE)minDE=c; if(c>maxDE)maxDE=c; }
    /* subset invariants */
    int bG=0,fG=0,treeG=0,pathG=0;
    for(U S=1;S<=full;S++){
      int k=pc(S);
      int e=edges(S);
      int comps=ncomp(S);
      int acyclic=(e==k-comps);
      if(acyclic){ if(k>fG) fG=k; if(comps==1){ if(k>treeG) treeG=k;
          int maxdeg=0; for(U m=S;m;m&=m-1){ int d=pc(adj[__builtin_ctz(m)]&S); if(d>maxdeg)maxdeg=d; }
          if(maxdeg<=2 && k>pathG) pathG=k; } }
      if(k>bG && bipartite(S)) bG=k;
    }
    /* Hamiltonian path DP */
    for(U S=0;S<=full;S++) for(int v=0;v<n;v++) hp[S][v]=0;
    for(int v=0;v<n;v++) hp[1u<<v][v]=1;
    for(U S=1;S<=full;S++) for(int v=0;v<n;v++) if(hp[S][v]){ U nb=adj[v]&~S; for(U m=nb;m;m&=m-1){ int u=__builtin_ctz(m); hp[S|(1u<<u)][u]=1; } }
    int hamS[1<<12]; for(U S=1;S<=full;S++){ int h=0; for(int v=0;v<n;v++) if(hp[S][v]){h=1;break;} hamS[S]=h; }
    int hasHamPath=hamS[full];
    /* path cover number */
    pcov[0]=0;
    for(U S=1;S<=full;S++){ int low=S&(-S); int best=99; U rest=S&~low;
      for(U sub=rest;;sub=(sub-1)&rest){ U T=sub|low; if(hamS[T]){ int c=1+pcov[S&~T]; if(c<best) best=c; } if(sub==0) break; }
      pcov[S]=best; }
    int pG=pcov[full];
    /* residue (Havel-Hakimi) */
    int deg[16]; for(int v=0;v<n;v++) deg[v]=pc(adj[v]);
    int seq[16]; for(int i=0;i<n;i++) seq[i]=deg[i];
    int len=n, residue=0;
    while(1){ /* sort descending */
      for(int i=0;i<len;i++) for(int j=i+1;j<len;j++) if(seq[j]>seq[i]){int t=seq[i];seq[i]=seq[j];seq[j]=t;}
      if(len==0){ residue=0; break; }
      if(seq[0]==0){ residue=len; break; }
      int d=seq[0]; for(int i=1;i<=d && i<len;i++) seq[i]--; for(int i=0;i<len-1;i++) seq[i]=seq[i+1]; len--; }
    /* second smallest degree */
    int sd[16]; for(int i=0;i<n;i++) sd[i]=deg[i]; for(int i=0;i<n;i++) for(int j=i+1;j<n;j++) if(sd[j]<sd[i]){int t=sd[i];sd[i]=sd[j];sd[j]=t;}
    int sigma=sd[1];
    /* C4 (not necessarily induced): two distinct vertices with >=2 common neighbours */
    int hasC4=0; for(int a=0;a<n && !hasC4;a++) for(int c=a+1;c<n;c++) if(pc(adj[a]&adj[c])>=2){hasC4=1;break;}
    /* vertex connectivity */
    int kappa=n-1;
    for(U S=0;S<full;S++){ int k=pc(S); if(k>=kappa) continue; U R=full&~S; if(!connected(R)) kappa=k; }
    /* --- conjectures --- */
    if(argc>1){ printf("%s n=%d b=%d f=%d tree=%d path=%d p=%d res=%d diam=%d rad=%d sumEcc=%d maxL=%d sumL=%d alpha=%d minDE=%d maxDE=%d sigma=%d kappa=%d C4=%d ham=%d\n",line,n,bG,fG,treeG,pathG,pG,residue,diam,rad,sumEcc,maxL,sumL,alpha,minDE,maxDE,sigma,kappa,hasC4,hasHamPath); }
    int bad[8]={0};
    /* 19: b >= floor(avgEcc + maxL) */
    { int rhs=(sumEcc + n*maxL)/n; if(bG<rhs) bad[0]=1; }
    /* 40: f >= ceil((p+b+1)/2) */
    { int rhs=(pG+bG+1+1)/2; if(fG<rhs) bad[1]=1; }
    /* 61: f >= residue + ceil(diam/3) */
    { int rhs=residue+(diam+2)/3; if(fG<rhs) bad[2]=1; }
    /* 72: tree >= ceil(avgEcc + maxL/3) = ceil((3 sumEcc + n maxL)/(3n)) */
    { int num=3*sumEcc+n*maxL, den=3*n; int rhs=(num+den-1)/den; if(treeG<rhs) bad[3]=1; }
    /* 96: alpha <= 1 + minDE*(maxL-1) */
    { if(alpha > 1 + minDE*(maxL-1)) bad[4]=1; }
    /* 133: path >= rad + (avgL)^{cC4}; cC4=0 if hasC4 */
    { if(hasC4){ if(pathG < rad+1) bad[5]=1; } else { if(n*pathG < n*rad + sumL) bad[5]=1; } }
    /* 189: maxDE <= 1+sigma => Ham path */
    { if(maxDE<=1+sigma && !hasHamPath) bad[6]=1; }
    /* 199: tree-2 <= kappa => Ham path */
    { if(treeG-2<=kappa && !hasHamPath) bad[7]=1; }
    for(int i=0;i<8;i++) if(bad[i]){ viol[i]++; if(printed[i]<3){ printed[i]++; printf("VIOL %d n=%d %s | b=%d f=%d tree=%d path=%d p=%d res=%d diam=%d rad=%d sumEcc=%d maxL=%d sumL=%d alpha=%d minDE=%d maxDE=%d sigma=%d kappa=%d C4=%d ham=%d\n",
        ids[i],n,line,bG,fG,treeG,pathG,pG,residue,diam,rad,sumEcc,maxL,sumL,alpha,minDE,maxDE,sigma,kappa,hasC4,hasHamPath); fflush(stdout);} }
  }
  printf("graphs=%lld",cnt); for(int i=0;i<8;i++) printf(" C%d:%lld",ids[i],viol[i]); printf("\n");
  return 0;
}
\end{lstlisting}
\subsection*{A.4 \texttt{l2.c}: inequality \eqref{eq:cand}}
\begin{lstlisting}[language=C]
/* Test (L2) f >= res + nu_ind and (L2') alpha2 >= res + nu_ind on connected graphs from stdin */
#include <stdio.h>
#include <stdlib.h>
#include <string.h>
typedef unsigned int U; static int n; static U adj[16];
static int pc(U x){return __builtin_popcount(x);}
static void parse(const char*s){ n=s[0]-63; memset(adj,0,sizeof adj); const char*p=s+1; int bits=0,val=0;
  for(int j=1;j<n;j++) for(int i=0;i<j;i++){ if(!bits){val=(*p++)-63;bits=6;} bits--; if((val>>bits)&1){adj[i]|=1u<<j;adj[j]|=1u<<i;} } }
static U reach(U mask,int s){U seen=1u<<s,fr=seen; while(fr){U nf=0; for(U f=fr;f;f&=f-1) nf|=adj[__builtin_ctz(f)]&mask; nf&=~seen; seen|=nf; fr=nf;} return seen;}
static int ncomp(U m){int c=0; while(m){U r=reach(m,__builtin_ctz(m)); m&=~r; c++;} return c;}
static int eu[64],ev[64],m_; static int best_;
static void rec(int start,U blocked,int sz){ if(sz>best_) best_=sz; for(int t=start;t<m_;t++){ U eb=(1u<<eu[t])|(1u<<ev[t]); if(eb&blocked) continue; rec(t+1,blocked|eb|adj[eu[t]]|adj[ev[t]],sz+1);} }
int main(){ char line[256]; long long cnt=0,b2=0,b2p=0; int sh=0;
  while(fgets(line,sizeof line,stdin)){ line[strcspn(line,"\n")]=0; parse(line); if(n<2) continue; cnt++;
    U full=(1u<<n)-1; int f=0,a2=0;
    for(U S=1;S<=full;S++){ int k=pc(S); int e=0,maxd=0; for(U x=S;x;x&=x-1){int d=pc(adj[__builtin_ctz(x)]&S); e+=d; if(d>maxd)maxd=d;} e/=2;
      if(maxd<=1 && k>a2) a2=k;
      if(k>f && e==k-ncomp(S)) f=k; }
    m_=0; for(int i=0;i<n;i++) for(int j=i+1;j<n;j++) if(adj[i]>>j&1){eu[m_]=i;ev[m_]=j;m_++;}
    best_=0; rec(0,0,0); int nu=best_;
    int seq[16],len=n,res=0; for(int i=0;i<n;i++) seq[i]=pc(adj[i]);
    while(1){ for(int i=0;i<len;i++) for(int j=i+1;j<len;j++) if(seq[j]>seq[i]){int t=seq[i];seq[i]=seq[j];seq[j]=t;}
      if(len==0){res=0;break;} if(seq[0]==0){res=len;break;} int d=seq[0]; for(int i=1;i<=d&&i<len;i++) seq[i]--; for(int i=0;i<len-1;i++) seq[i]=seq[i+1]; len--; }
    if(f<res+nu){ b2++; if(sh++<5){printf("L2 FAIL %s f=%d res=%d nu=%d\n",line,f,res,nu); fflush(stdout);} }
    if(a2<res+nu){ b2p++; }
  }
  printf("graphs=%lld L2fail=%lld L2'(alpha2)fail=%lld\n",cnt,b2,b2p);
}
\end{lstlisting}
\end{document}
